
\section{Discussion}

\paragraph{Computing $\occmeaselem$ in high dimensions:} 
Our early stopping method requires computing the occupancy measure $\occmeaselem$. Occupancy measures can be approximated via rollouts, though this approximation may be expensive and noisy.
Another option is to solve for $\occmeaselem = \occmeas{\policy}$ via  $\vec{\occmeaselem} = (I-\Pi \TMat)^{-1}\Pi \vec{\InitStateDistribution}$ where $\TMat$ is the transition matrix, $\InitStateDistribution$ is the initial state distribution, and $\Pi_{s, (s, a)} = \mathbb{P}(\pi(s) = a)$. This solution could be approximated in large environments. 

\paragraph{Approximating $\theta$:} Our early stopping method requires an upper bound $\theta$ on the angle between the true reward and the proxy reward. In practice, this should be seen as a measure of how accurate we believe the proxy to be. If the proxy reward is obtained through reward learning, then we may be able to estimate $\theta$ based on the learning algorithm, the amount of training data, and so on. 
Moreover, if we have a (potentially expensive) method to evaluate the true reward, such as expert judgement, then we can estimate $\theta$ directly (even in large environments). For details, see \citet{skalse2023starc}.

\paragraph{Key assumptions:} An important consideration when employing any optimisation algorithm is its behaviour when its key assumptions are not met. 
For our early stopping method, if the provided $\theta$ does not upper-bound the angle between the proxy and the true reward, then the learnt policy may, in the worst case, result in as much Goodharting as a policy produced by na\"ive optimisation.\footnote{However, it might still be possible to bound the worst-case performance further using the norm of the transition matrix (defining the geometry of the polytope $\Omega$). This will be an interesting topic for future work.} On the other hand, if the optimisation algorithm is not concave, then this can only cause the early-stopping procedure to stop at a sub-optimal point; Goodharting is still guaranteed to be avoided. This is also true if the upper bound $\theta$ is not tight.

\paragraph{Significance and Implications:} Our work has several direct implications.
In Section~\ref{sec:quantifying_goodharting}, we show that Goodharting occurs for a wide range of environments and reward functions. This means that we should expect to see Goodharting often when optimising for misspecified proxy rewards. 
In Section~\ref{sec:explain_goodhart}, we provide a mechanistic explanation for \emph{why} Goodharting occurs. We expect this to be helpful for further progress in the study of reward misspecification.
In Section~\ref{sec:preventing_goodhart}, we provide early stopping methods that provably avoid Goodharting, and show that these methods, in a certain sense, are worst-case optimal.
However, these methods can lead to less true reward than na\"ive optimisation,
This means that they are most applicable when it is essential to avoid Goodharting. 

\paragraph{Limitations and Future Work:}
We do not have a comprehensive understanding of the dynamics at play when a misspecified reward function is maximised, and our work does not exhaust this area of study.
An important question is what types of failure modes can occur in this setting, and how they may be detected and mitigated. Our work studies one important failure mode (i.e.\ Goodharting), but there may be other distinctive failure modes that could be described and studied as well. A related important question is precisely how a proxy reward $\ProxR$ may differ from the true reward $\TrueR$, before maximising $\ProxR$ might be bad according to $\TrueR$. There are several existing results pertaining to this question \citep{ng_policy_1999, gleave2021quantifying, skalse_defining_2022, skalse2022invariance}, but there is at the moment no comprehensive answer. Another interesting direction is to use our results to develop policy optimisation algorithms that collect more data about the reward function over time, as this information is needed. We discuss this direction in Appendix~\ref{appendix:iterative_improvement}. Finally, it would be interesting to try to find principled relaxations of the methods in Section~\ref{sec:preventing_goodhart}, that attain better practical performance while retaining desirable theoretical guarantees. 
